% ═══════════════════════════════════════════════════════════════
% LH-01b: Lehrerhinweise - Sich vorstellen + Verb ser
% Abgeleitet aus LE-01b (Verlauf, Differenzierung, Fehler)
% ═══════════════════════════════════════════════════════════════

\documentclass[11pt, a4paper]{article}

\usepackage[utf8]{inputenc}
\usepackage[T1]{fontenc}
\usepackage[ngerman]{babel}
\usepackage{helvet}
\renewcommand{\familydefault}{\sfdefault}

\usepackage[margin=2cm]{geometry}
\usepackage{enumitem}
\usepackage{tabularx}
\usepackage{booktabs}
\usepackage{longtable}

\usepackage{xcolor}
\usepackage{amssymb}
\usepackage{tcolorbox}

\usepackage{fancyhdr}
\usepackage{lastpage}

% Farben
\definecolor{spanisch}{HTML}{c41e3a}
\definecolor{grau}{HTML}{888888}
\definecolor{hellgrau}{HTML}{f5f5f5}
\definecolor{basis}{HTML}{2a7a4b}
\definecolor{standard}{HTML}{2c5aa0}
\definecolor{erweiterung}{HTML}{7b1fa2}
\definecolor{loesung}{HTML}{c62828}

\newcommand{\fachfarbe}{spanisch}

% Variablen
\newcommand{\thema}{Sich vorstellen + Verb ser}
\newcommand{\sequenz}{01}
\newcommand{\block}{b}
\newcommand{\fach}{Spanisch}
\newcommand{\klasse}{7}
\newcommand{\dauer}{90 Min. (Doppelstunde)}

% Kopf-/Fußzeile
\pagestyle{fancy}
\fancyhf{}
\renewcommand{\headrulewidth}{0.4pt}
\renewcommand{\footrulewidth}{0.4pt}
\lhead{\fach{} -- Klasse \klasse}
\chead{\textbf{LH-\sequenz\block{} (Lehrerhinweise)}}
\rhead{}
\lfoot{}
\rfoot{Seite \thepage\ von \pageref{LastPage}}

% Differenzierungsmarker
\newcommand{\B}{\textcolor{basis}{\textbf{[B]}}}
\newcommand{\Std}{\textcolor{standard}{\textbf{[S]}}}
\newcommand{\E}{\textcolor{erweiterung}{\textbf{[E]}}}
\newcommand{\loesung}[1]{\textcolor{loesung}{\textbf{#1}}}
\newcommand{\es}[1]{\textcolor{\fachfarbe}{\textbf{#1}}}

% Boxen
\newtcolorbox{kurzplan}{
    colback=hellgrau,
    colframe=\fachfarbe,
    boxrule=1pt,
    arc=0pt,
    fonttitle=\bfseries,
    title=Kurzplan
}

\newtcolorbox{differenzierung}{
    colback=white,
    colframe=grau,
    boxrule=0.5pt,
    arc=0pt,
    fonttitle=\bfseries,
    title=Differenzierung
}

\newtcolorbox{fehlerbox}{
    colback=white,
    colframe=loesung,
    boxrule=0.5pt,
    arc=0pt,
    fonttitle=\bfseries,
    title=Häufige Fehler
}

\newtcolorbox{materialbox}{
    colback=white,
    colframe=\fachfarbe,
    boxrule=0.5pt,
    arc=0pt,
    fonttitle=\bfseries,
    title=Material-Checkliste
}

% ═══════════════════════════════════════════════════════════════
\begin{document}

\begin{center}
    {\Large\bfseries\textcolor{\fachfarbe}{Lehrerhinweise: \thema}}\\[0.3em]
    {\normalsize LH-\sequenz\block{} | \dauer{} | Std. 3-4 von 8}
\end{center}

\vspace{10pt}

% ═══════════════════════════════════════════════════════════════
% 1. KURZPLAN (aus LE-01b Verlauf)
% ═══════════════════════════════════════════════════════════════

\section*{1. Stundenverlauf (Kurzplan)}

\textbf{1. Stunde (45 Min.)}

\begin{tabularx}{\textwidth}{|c|l|X|l|}
\hline
\textbf{Zeit} & \textbf{Phase} & \textbf{Inhalt} & \textbf{Material} \\
\hline
0--5 & Einstieg & Audio: Vorstellung anhören & Audio \\
\hline
5--10 & Aktivierung & Redemittel sammeln (Name, Ort) & Tafel \\
\hline
10--20 & Input & Verb \textit{ser} einführen, Formen ableiten & Tafel, AB \\
\hline
20--30 & Übung 1 & Chorisches Sprechen: Konjugation & -- \\
\hline
30--40 & Übung 2 & PA: Minidialoge mit Karten & Karten \\
\hline
40--45 & Sicherung & Paare präsentieren & -- \\
\hline
\end{tabularx}

\vspace{1em}

\textbf{2. Stunde (45 Min.)}

\begin{tabularx}{\textwidth}{|c|l|X|l|}
\hline
\textbf{Zeit} & \textbf{Phase} & \textbf{Inhalt} & \textbf{Material} \\
\hline
0--5 & Aktivierung & Quiz: ser-Formen abfragen & -- \\
\hline
5--15 & Vertiefung & EA: Lückentext AB-01b & AB-01b \\
\hline
15--25 & Anwendung & PA/GA: Rollenspiel ``Neue kennenlernen'' & Namensschilder \\
\hline
25--35 & Erweiterung & Kulturvergleich tú/usted (\E{}) & Tafel \\
\hline
35--40 & Sicherung & Merkkasten gemeinsam ausfüllen & AB-01b \\
\hline
40--45 & Abschluss & Zusammenfassung, HA & -- \\
\hline
\end{tabularx}

% ═══════════════════════════════════════════════════════════════
% 2. DIFFERENZIERUNG (aus LE-01b)
% ═══════════════════════════════════════════════════════════════

\section*{2. Differenzierung}

\begin{differenzierung}
\textbf{Legende:} \B{} Basis (BR) | \Std{} Standard (MR) | \E{} Erweiterung (GY)

\vspace{0.5em}
\textbf{WICHTIG:} Diese Marker sind NUR für die Lehrkraft!
\end{differenzierung}

\vspace{1em}

\begin{tabularx}{\textwidth}{|l|X|X|X|}
\hline
& \B{} \textbf{Basis} & \Std{} \textbf{Standard} & \E{} \textbf{Erweiterung} \\
\hline
\textbf{Aufg. 1} & Zuordnung mit Hilfe (Tabelle sichtbar) & Zuordnung selbstständig & + Sätze bilden \\
\hline
\textbf{Aufg. 2} & Wortbank mit allen Formen & Wortbank reduziert & Ohne Wortbank \\
\hline
\textbf{Aufg. 3} & Dialog mit Vorlage & Dialog mit Stichworten & Freier Dialog \\
\hline
\textbf{Aufg. 4} & Optional / verkürzt (4 Zeilen) & Vollständig (6 Zeilen) & + Kulturvergleich tú/usted \\
\hline
\end{tabularx}

\vspace{1em}

\textbf{Konkrete Hinweise:}
\begin{itemize}
    \item \B{} LRS-SuS: Zeitzugabe (+10 Min.), ser-Tabelle als Hilfe erlauben
    \item \B{} DaZ-SuS: Deutsch-Spanisch Wortschatz vorab besprechen
    \item \Std{} Standardgruppe: Partnerarbeit mit wechselnden Partnern
    \item \E{} Leistungsstarke: Kulturvergleich, Expertenrolle, Zusatzaufgabe 4
\end{itemize}

% ═══════════════════════════════════════════════════════════════
% 3. HÄUFIGE FEHLER (aus LE-01b Abschnitt 7)
% ═══════════════════════════════════════════════════════════════

\section*{3. Häufige Fehler}

\begin{fehlerbox}
\begin{tabularx}{\textwidth}{|l|X|X|}
\hline
\textbf{Fehler} & \textbf{Beispiel} & \textbf{Reaktion} \\
\hline
Aussprache \textit{eres} & *[eres] statt ['eres] & Vorsprechen, chorisch wiederholen \\
\hline
Verwechslung soy/es & *``Él soy de...'' & Kontrastive Übung: 1. vs. 3. Person \\
\hline
Dopplung durch Interferenz & *``Yo soy soy María'' & Bewusstmachung: Ein Verb reicht! \\
\hline
Pronomen vergessen & *``Soy María'' ohne Kontext & Im Spanischen optional, aber zur Verdeutlichung hilfreich \\
\hline
\end{tabularx}
\end{fehlerbox}

% ═══════════════════════════════════════════════════════════════
% 4. MATERIAL-CHECKLISTE
% ═══════════════════════════════════════════════════════════════

\section*{4. Material-Checkliste}

\begin{materialbox}
\begin{itemize}[label=$\square$]
    \item AB-01b.pdf (24 Kopien)
    \item ML-01b.pdf (1× für Lehrkraft)
    \item Lehrwerk: Qué pasa 1, S. 10--15
    \item Audio Track 5-6 (CD1 / USB)
    \item Redemittel-Karten (Kopiervorlage, 12 Sets)
    \item Namensschilder für Rollenspiel
    \item Optional: LP-01b.html (Tablets/PCs)
\end{itemize}
\end{materialbox}

% ═══════════════════════════════════════════════════════════════
% 5. LÖSUNGEN (aus LE-01b Abschnitt 9)
% ═══════════════════════════════════════════════════════════════

\section*{5. Lösungen AB-01b}

\textbf{Merkkasten:}
yo $\rightarrow$ soy, tú $\rightarrow$ eres, él/ella $\rightarrow$ es, nosotros $\rightarrow$ somos

\vspace{0.5em}

\textbf{Aufgabe 1:} (je 1P)
\begin{itemize}
    \item yo $\rightarrow$ soy
    \item tú $\rightarrow$ eres
    \item él/ella $\rightarrow$ es
    \item nosotros $\rightarrow$ somos
\end{itemize}

\textbf{Aufgabe 2:} (je 1P)
\begin{enumerate}[label=\alph*)]
    \item \loesung{Soy}
    \item \loesung{Soy}
    \item \loesung{es}
    \item \loesung{somos}
\end{enumerate}

\textbf{Aufgabe 3:} (je 1P, max. 6P)
\begin{itemize}
    \item A: ¿Cómo te llamas?
    \item B: Me llamo [Name].
    \item A: ¿De dónde eres?
    \item B: Soy de [Ort/Land].
    \item B: ¡Igualmente! / ¡Mucho gusto!
\end{itemize}

\textbf{Aufgabe 4:} (6P)
\begin{itemize}
    \item Begrüßung: 1P
    \item Name: 1P
    \item Herkunft: 2P
    \item Verabschiedung: 1P
    \item Sprachliche Korrektheit: 1P
    \item Bonus (tú/usted): +1P
\end{itemize}

% ═══════════════════════════════════════════════════════════════
% 6. HAUSAUFGABE
% ═══════════════════════════════════════════════════════════════

\section*{6. Hausaufgabe}

\begin{itemize}
    \item Lehrwerk S. 12, Übung 3 (ser einsetzen)
    \item Vokabeln lernen: ser-Formen, Redemittel
    \item Optional: Lernpfad LP-01b.html bearbeiten
\end{itemize}

% ═══════════════════════════════════════════════════════════════
% 7. REFLEXIONSNOTIZEN
% ═══════════════════════════════════════════════════════════════

\section*{7. Reflexionsnotizen (nach Durchführung)}

\begin{tabularx}{\textwidth}{|l|X|}
\hline
\textbf{Was lief gut?} & \\[2em]
\hline
\textbf{Was lief nicht?} & \\[2em]
\hline
\textbf{Zeitmanagement} & \\[2em]
\hline
\textbf{Für nächstes Mal} & \\[2em]
\hline
\end{tabularx}

\end{document}
